\section{符号表与约定}
\label{app:symbols}
\label{app:conventions}

本附录是全书符号的\textbf{唯一权威}。正文任何一章首次使用符号时都给出定义；若与
本附录冲突，以本附录为准。引用河野、Moore、李思、上野四家的公式时，必须先在下表
查出其归一化，换算成本文约定并在页边注\texttt{[换算见 appA]}，不得直接照抄。

\subsection{全局符号表}

\begin{center}
\small
\begin{longtable}{@{}>{\raggedright\arraybackslash}p{2.3cm}>{\raggedright\arraybackslash}p{9.4cm}>{\raggedright\arraybackslash}p{2.5cm}@{}}
\toprule
符号 & 定义与约定 & 备注 \\
\midrule
\endfirsthead
\toprule
符号 & 定义与约定 & 备注 \\
\midrule
\endhead
$M$ & 定向、闭（通常无边）三维流形 & \src{353} \\
$\Sigma$ & 定向曲面，亏格 $g$、$n$ 个标记点；$M=\Sigma\times\mathbb R$ 时它是空间切片 & \src{366} \\
$G$ & 紧、单连通、单 Lie 群；$d=\dim G$；$h^\vee$ 对偶 Coxeter 数；$\mathfrak g=\mathrm{Lie}\,G$ & \src{353} \\
$\Tr$ & 取 $G$ 的\textbf{定义表示}（$SU(N)$ 即 $N$ 维表示）的迹；长根平方 $=2$ & \src{353} 脚注 1 \\
$A$ & $G$-联络，$\mathfrak g$-值 1-形式，$A=A_i\dd x^i$；取反厄米约定 & \src{353} \\
$F$ & 曲率 2-形式，$F=\dd A+A\wedge A$；分量 $F_{ij}=\partial_iA_j-\partial_jA_i+[A_i,A_j]$ & \src{353} \\
$D_i$ & 协变导数 $D_i=\partial_i+[A_i,\cdot\,]$；无穷小规范变换 $A_i\to A_i-D_i\varepsilon$ & \src{353} \\
$L$ & 陈--西蒙斯作用量 $\frac{k}{4\pi}\int_M\Tr(A\wedge\dd A+\frac23A^3)$ & \src{354} \\
$k$ & 陈--西蒙斯能级，$k\in\mathbb Z$（由大规范变换单值性定出） & \src{354} \\
$W_R(C)$ & 沿闭曲线 $C$、表示 $R$ 的 Wilson 线（定义式 (1.5)） & \src{355} \\
$\mathcal A$ & $\Sigma$ 上所有 $G$-联络的仿射空间 & \src{366} \\
$\mathcal M$ & 平坦联络模空间 $\Hom(\pi_1\Sigma,G)/G$（辛商） & \src{367} \\
$\mathcal M_J$ & 选定复结构 $J$ 后，$\mathcal M$ 作为全纯（Kähler）流形的版本；与 $\mathcal M$ 不混用 & \src{368} \\
$\omega$ & $\mathcal M$ 上的辛形式 $\frac{k}{4\pi}\int_\Sigma\Tr(\delta A\wedge\delta A)$ & \src{367} \\
$\mathcal L$ & $\mathcal M_J$ 上 $\bar\partial$ 族的行列式线丛；带 Quillen 度规 & \src{369} \\
$\mathcal H_\Sigma^{(J)}$ & 量子 Hilbert 空间 $=H^0(\mathcal M_J,\mathcal L^{\otimes k})$；上标 $(J)$ 表示依赖复结构 & \src{369} \\
$\lambda$ & 最高权；$\rho=\frac12\sum_{\alpha>0}\alpha$；共形权 $h_\lambda=\frac{(\lambda,\lambda+2\rho)}{2(k+h^\vee)}$ & \src{369} \\
$q$ & $SU(N)$：$q=e^{2\pi\ii/(N+k)}$；一般 $G$：$q=e^{2\pi\ii/(k+h^\vee)}$ & \src{382} \\
$S_{ij}$ & 模变换 $\tau\to-1/\tau$ 的矩阵；$SU(2)$ 显式见 \eqref{eq:appA-S} & \src{390} \\
$N_{ijk}$ & 融合系数；$S^2\times S^1$ 上三点电荷 Hilbert 空间的维数 & \src{388} \\
$K_i{}^j$ & 手术矩阵，把 $M$ 上的配分函数表成 $\widetilde M$ 上的配分函数 & \src{385} \\
$Z(M)$ & 配分函数（路径积分）；$Z(S^3)=S_{0,0}$ & \src{389--390} \\
\bottomrule
\end{longtable}
\end{center}

\subsection{术语中英对照}

\begin{center}
\small
\begin{longtable}{@{}p{3.1cm}p{5.1cm}p{5.9cm}@{}}
\toprule
中文 & 英文 & 备注 \\
\midrule
\endfirsthead
\toprule
中文 & 英文 & 备注 \\
\midrule
\endhead
陈--西蒙斯作用量 & Chern--Simons action & 本文 $L$；河野记 $\mathrm{CS}$ \\
威尔逊线 & Wilson line & 沿链环的迹 Holonomy \\
穿刺点 / 标记点 & puncture / marked point & 带电荷的静态源位置 \\
共形块 & conformal block & $\mathcal H_\Sigma$ 的全纯解释 \\
融合规则 & fusion rule & $N_{ijk}\neq0$ 的条件 \\
射影平坦联络 & projectively flat connection & $\mathcal M$ 上的 $(1,0)$ 联络 \\
极化 & polarization & 选全纯截面子空间 \\
预量子化 & prequantization & 线丛+联络，曲率 $=\omega/\hbar$ \\
全纯截面 & holomorphic section & $H^0(\mathcal M_J,\mathcal L^{\otimes k})$ \\
行列式线丛 & determinant line bundle & $\det\ker\bar\partial\otimes\det(\mathrm{coker}\,\bar\partial)^*$ \\
半形式修正 & half-form (metaplectic) correction & 本文未引入，见 \S\ref{app:errata} \\
标架反常 & framing anomaly & 相位 $e^{2\pi\ii c/24}$ \\
手术 & surgery & 挖去实心环面再粘回 \\
Dehn 手术 & Dehn surgery & 河野用它定义 $Z_k(M)$ \\
链环 & link & 多个纽结的并 \\
辫群 & braid group & Jones 多项式的代数来源 \\
模变换 & modular transformation & $S:\tau\to-1/\tau$，$T:\tau\to\tau+1$ \\
中心荷 & central charge & $c=\dfrac{k\dim G}{k+h^\vee}$ \\
可积表示 & integrable representation & 最高权 $\lambda$ 满足 $(\theta,\lambda)\le k$ \\
对偶 Coxeter 数 & dual Coxeter number & $h^\vee$；$SU(N)$ 为 $N$ \\
顶点算子 & vertex operator & 共形场的算子实现 \\
算符乘积展开 & operator product expansion (OPE) & 顶点代数的乘法 \\
拓扑量子场论 & topological quantum field theory (TQFT) & Atiyah--Segal 公理 \\
\bottomrule
\end{longtable}
\end{center}

\subsection{本文的出发点}

本文只有一条归一化需要记住：作用量取
\begin{equation}
L=\frac{k}{4\pi}\int_M\Tr\Bigl(A\wedge\dd A+\frac23A^3\Bigr),\qquad k\in\mathbb Z,
\label{eq:appA-cs}
\end{equation}
$\Tr$ 取定义表示、长根平方 $=2$。$k$ 的整性来自大规范变换下单值性
\src{354}；它同时固定了 $\mathcal M$ 上辛形式的整性归一化
$\omega=\frac{k}{4\pi}\int_\Sigma\Tr(\delta A\wedge\delta A)$ \src{367}。
式 \eqref{eq:appA-cs} 与 \eqref{eq:cs} 是同一个式子，后者是正文首次引入时的编号。

\subsection{与四本参考书的换算}

\subsubsection*{陈--西蒙斯作用量：$L_{\text{本文}}=2\pi k\,\mathrm{CS}_{\text{河野}}$}

河野俊丈的 CS 泛函（中译 (2.33)，p.~71）取
$\mathrm{CS}_{\text{河野}}(A)=\dfrac{1}{8\pi^2}\displaystyle\int_M\Tr\bigl(A\dd A+\tfrac23A^3\bigr)$，
即系数 $1/(8\pi^2)$；本文系数为 $k/(4\pi)$。两者之比是纯数字：
\begin{equation}
\frac{k}{4\pi}\Big/\frac{1}{8\pi^2}=\frac{k}{4\pi}\cdot 8\pi^2=2\pi k
\quad\Longrightarrow\quad
L_{\text{本文}}=2\pi k\,\mathrm{CS}_{\text{河野}}.
\label{eq:appA-cs-rel}
\end{equation}
于是路径积分的权重完全一致：
\begin{equation}
e^{\ii L}=e^{2\pi\ii k\,\mathrm{CS}_{\text{河野}}}.
\end{equation}
河野证明 $\mathrm{CS}_{\text{河野}}(g^*A)=\mathrm{CS}_{\text{河野}}(A)-$（整数）
（中译 (2.36)，p.~71；映射度），故 $e^{2\pi\ii k\mathrm{CS}}$ 在大规范变换下单值当且
仅当 $k\in\mathbb Z$——这正是本文 \src{354} 的整性条件。也就是说：\textbf{两本书的 $k$
是同一个整数能级}，差别只在作用量的整体常数 $2\pi k$。

辛形式也要对齐。本文 $\omega_{\text{本文}}=\frac{k}{4\pi}\int_\Sigma\Tr(\delta A\wedge\delta A)$；
河野（中译 (2.37)，p.~72）取 $\omega_{\text{河野}}=-\frac{1}{8\pi^2}\int_\Sigma\Tr(\alpha\wedge\beta)$，
能级 $k$ 时乘以 $k$。代入得
\begin{equation}
\omega_{\text{本文}}=\frac{k}{4\pi}\int\Tr=-\frac{k}{4\pi}\Bigl(-\frac{1}{8\pi^2}\Big/\frac{1}{8\pi^2}\Bigr)
\int\Tr=-2\pi k\,\omega_{\text{河野}}.
\label{eq:appA-omega}
\end{equation}
（化简：$-2\pi k\cdot\bigl(-\frac{1}{8\pi^2}\bigr)=+\frac{k}{4\pi}$，故
$\omega_{\text{本文}}=-2\pi k\,\omega_{\text{河野}}$。）注意这里同时出现整体符号与
$2\pi k$ 因子：3-形式积分同号，但“边界辛形式”的符号约定相反（河野从 $\delta\mathrm{CS}$
的边界项取负号）。比较泊松括号或量子化条件时，必须用同一个 $\omega$，否则符号会错。

\subsubsection*{能级：$k$、$\ell$、$\kappa$ 不是同一个量}

\begin{itemize}[leftmargin=1.6em]
\item \textbf{本文}：$k$ 是陈--西蒙斯能级，可积表示满足 $(\theta,\lambda)\le k$
（$SU(2)$ 即 $2j\le k$）；公式里的 $k+h^\vee$ 是\textbf{移位能级}，不是另一个 level。
\item \textbf{河野}：$k$ 同义（可积最高权 $\lambda\le k$）；$SU(2)$ 的 $k+2$ 就是
$k+h^\vee$（$h^\vee_{SU(2)}=2$）。
\item \textbf{Moore}：非阿贝尔情形用 $k$；阿贝尔 $U(1)$ 情形用 $\kappa$，作用量
$(\kappa/2\pi)\int A\wedge\dd A$，$\kappa\in\mathbb Z$；两者换算 $k=-2\kappa$
（$U(1)$ 的 $2\kappa$ 维 Hilbert 空间对应非阿贝尔的能级 $k$）。\textbf{此 $\kappa$ 与上野的
$\kappa=\ell+g^*$ 无关}，不要混。
\item \textbf{李思}：没有“能级”，只有线丛张量幂 $L=L_0^{\otimes k}$ 的整数 $k$；在
$\hbar=1/(2\pi k)$ 下它就是本文的能级（见下）。
\item \textbf{上野}：书里的 level 记作 $\ell$（$P_\ell=\{\lambda:0\le(\theta,\lambda)\le\ell\}$，
p.~39）；涉及中心荷与共形权时用 $\kappa=\ell+g^*$，$g^*=h^\vee$（p.~125 批注：
“反常移位 $\ell\to\ell+g^*$”）。这就是半形式修正，本文略去（见 \S\ref{app:errata}）。
\end{itemize}

\subsubsection*{Planck 常数：$\hbar=1/(2\pi k)$}

李思的预量子化（定义 5.4.3，p.~303--304）要求
\begin{equation}
\frac{1}{2\pi\hbar}\int_\Sigma\omega\in\mathbb Z,\qquad
F_\nabla=-\frac{\ii}{\hbar}\omega .
\label{eq:appA-preq}
\end{equation}
例 5.4.11（p.~308）取 $L=L_0^{\otimes k}$，曲率
$F_\nabla=kF_{\nabla0}=-\ii(2\pi k)\omega$，并明说：若认同
$\hbar=1/(2\pi k)$，则 $(L,\nabla)$ 给出预量子化线丛。把它与本文对齐：本文
$k=1$ 时 $c_1(\mathcal L)=[\omega]$（\src{369}），故 $c_1(\mathcal L^{\otimes k})=k[\omega]$，
即 $\int_{S^2}\omega=1$（对生成元取归一化）。代入 \eqref{eq:appA-preq}：
$\frac{1}{2\pi\hbar}\cdot k\in\mathbb Z$ 的最小选择给出
\begin{equation}
\hbar=\frac{1}{2\pi k}\qquad\Longleftrightarrow\qquad k=\frac{1}{2\pi\hbar}.
\end{equation}
\textbf{陷阱}：数学教材写 $c_1=[\ii F/2\pi]$，本文写“$c_1(\mathcal L)=[\omega]$、曲率
$-\ii\omega$”，两者相差一个 $2\pi$；因此不能把 Witten 的“曲率 $=-\ii\omega$”直接
代入 $F=-(\ii/\hbar)\omega$ 去解 $\hbar$（会得到 $\hbar=1$）。以上整性条件为准。
若采用 $\int_{S^2}\omega=2\pi$ 的 Fubini--Study 归一，则 $\hbar=1/k$——不同教材的 $2\pi$
差别全在这里。

\subsubsection*{$q$ 的定义}

本文 $SU(N)$ 的 $q=e^{2\pi\ii/(N+k)}$ \src{382}，一般群
$q=e^{2\pi\ii/(k+h^\vee)}$。河野的模变换用 $1/(k+2)$（$SU(2)$，中译 (1.42)），
与 $h^\vee_{SU(2)}=2$ 一致；Moore 用 $q=\exp(2\pi\ii/(k+h))$（$h$ 即 $h^\vee$），
完全同式；上野的 $S$ 矩阵含 $e^{-2\pi\ii(\lambda+\rho)/(\ell+g^*)}$，故
$q=e^{2\pi\ii/(\ell+g^*)}$，$\ell$ 是 level。四家只差把 $h^\vee$ 写成 $2$、$h$、$g^*$
还是 $\ell$ 的记号。

\subsubsection*{$S$ 矩阵}

本文（$SU(2)$，\src{390}）：
\begin{equation}
S_{mn}=\sqrt{\frac{2}{k+2}}\;\sin\frac{(m+1)(n+1)\pi}{k+2},
\qquad m,n=0,\dots,k .
\label{eq:appA-S}
\end{equation}
河野 (1.42)（p.~36）与此逐字相同（把 $m+1$ 写作 $\lambda+1$）。上野
(5.29)--(5.30)（p.~111--115）写成“特征标”形式
\begin{equation}
S_{\lambda\mu}=S_{\lambda0}\,\chi_\mu\!\left(e^{-2\pi\ii(\lambda+\rho)/(\ell+g^*)}\right),
\qquad
S_{\lambda0}=\Bigl|\frac{P}{(\ell+g^*)Q^\vee}\Bigr|^{-1/2}
\prod_{\alpha>0}2\sin\frac{\pi(\lambda+\rho,\alpha)}{\ell+g^*} .
\label{eq:appA-Sueno}
\end{equation}
$SU(2)$ 验证：$|P/Q^\vee|=2$，故
$S_{\lambda0}^{\text{上野}}=\sqrt{\tfrac{k+2}{2}}\cdot 2\sin\frac{\pi(\lambda+1)}{k+2}
=\sqrt{2(k+2)}\sin\frac{\pi(\lambda+1)}{k+2}=(k+2)\,S_{0\lambda}^{\text{本文}}$。
即
\begin{equation}
S^{\text{上野}}=(k+h^\vee)\,S^{\text{本文}}\quad(\text{对 }SU(2)\text{ 为 }(k+2)\text{ 倍}).
\label{eq:appA-Srel}
\end{equation}
Verlinde 公式
$N_{ijk}=\sum_m S_{im}S_{jm}S_{km}/S_{0m}$ 的分子、分母同乘一个常数不变，
故 \eqref{eq:appA-Srel} 的整体因子不影响任何物理结论；但单独引用某个矩阵元时必须
乘回 $(k+h^\vee)$。Schottenloher 全书没有 $S$ 矩阵，只给维数公式 (11.2)（p.~177），
不能当 $S$ 矩阵出处。

\subsubsection*{数值检验：$SU(2)$、$k=1$}

取 $SU(2)$、$k=1$（$h^\vee=2$，$q=e^{2\pi\ii/3}$）。本文 \eqref{eq:appA-S} 给出
\begin{equation}
S^{\text{本文}}=\sqrt{\frac23}
\begin{pmatrix}
\sin\frac{\pi}{3} & \sin\frac{2\pi}{3}\\[2pt]
\sin\frac{2\pi}{3} & \sin\frac{4\pi}{3}
\end{pmatrix}
=\frac{1}{\sqrt2}
\begin{pmatrix}1&1\\1&-1\end{pmatrix},
\end{equation}
即 $S_{00}=S_{01}=1/\sqrt2\approx0.7071$。上野 \eqref{eq:appA-Sueno} 的
$S_{\lambda0}=\sqrt{2(k+2)}\sin\frac{\pi(\lambda+1)}{k+2}=\sqrt6\sin\frac{(\lambda+1)\pi}{3}$，
$\lambda=0,1$ 都等于 $\sqrt6\cdot\frac{\sqrt3}{2}=\frac{3}{\sqrt2}\approx2.1213$。比值
\begin{equation}
\frac{S_{\lambda0}^{\text{上野}}}{S_{0\lambda}^{\text{本文}}}=3=k+2=k+h^\vee,
\end{equation}
与 \eqref{eq:appA-Srel} 一致。再用 Verlinde 公式核一个融合系数：$k=1$ 时
$j\le1/2$，$N_{1,1,0}$ 应为 $1$（$1/2\otimes1/2$ 只含自旋 $0$），而
\begin{equation}
N_{110}=\sum_{\alpha=0}^{1}\frac{S_{1\alpha}S_{1\alpha}S_{0\alpha}}{S_{0\alpha}}
=\sum_{\alpha=0}^{1}S_{1\alpha}^2=\frac12+\frac12=1 .
\end{equation}
三项核对（$S$ 的显式、上野的整体因子、Verlinde 公式）都通过，说明上面的换算没有
漏掉 $2\pi$ 或 $k+h^\vee$。

\subsection{速查表}

\begin{center}
\footnotesize
\begin{longtable}{@{}p{1.8cm}p{2.5cm}p{2.4cm}p{2.4cm}p{3.4cm}@{}}
\toprule
量 & 本文 & 河野 & Moore & 李思 / 上野 \\
\midrule
\endfirsthead
\toprule
量 & 本文 & 河野 & Moore & 李思 / 上野 \\
\midrule
\endhead
作用量 & $\frac{k}{4\pi}\int\Tr$ & $\frac{1}{8\pi^2}\int\Tr$ & $\frac{\kappa}{2\pi}\int A\dd A$（$U(1)$） & 李思：无 CS \\
换算 & — & $L=2\pi k\,\mathrm{CS}$ & $k=-2\kappa$ & — \\
能级 & $k$ & $k$ & $k$ / $\kappa$ & 李思：$k$（张量幂）\newline 上野：$\ell$ \\
移位 & $k+h^\vee$ & $k+2$ & $k+h$ & 上野：$\ell+g^*$，$g^*=h^\vee$ \\
$q$ & $e^{2\pi\ii/(N+k)}$ & $e^{2\pi\ii/(k+2)}$ & $e^{2\pi\ii/(k+h)}$ & 上野：$e^{2\pi\ii/(\ell+g^*)}$ \\
$S$ & 原文 (4.39) & (1.42) 同式 & — & 上野 (5.29)，$\times(k+h^\vee)$ \\
$\hbar$ & — & — & — & 李思：$\hbar=1/(2\pi k)$ \\
\bottomrule
\end{longtable}
\end{center}

\subsection{六个常见陷阱}

\begin{enumerate}[leftmargin=1.8em]
\item \textbf{$2\pi$ 账}：作用量的 $2\pi k$、辛形式的 $-2\pi k$、$c_1$ 定义里的 $2\pi$、
$\hbar=1/(2\pi k)$——四处必须同时对齐，只对一个会得到“$\hbar=1$”或“$k$ 差 $2\pi$”。
\item \textbf{level 与移位 level}：$k$ 是能级；$k+h^\vee$ 只出现在 $c$、$h_\lambda$、$q$
与上野的 $\ell+g^*$ 中。不要把“$2j\le k$”写成“$2j\le k+h^\vee$”。
\item \textbf{半形式}：本文没有 $K^{1/2}$ 修正，$k\to k+h^\vee$ 的“规范”说法来自译本
批注 \texttt{witten.tex:1206--1208} [译注]，不是 Witten 原文。
\item \textbf{符号}：$A_i\to A_i-D_i\varepsilon$（$D_i=\partial_i+[A_i,\cdot]$）是本文
约定 \src{353}；河野在同一处可能写成 $+$。引河野的 (3.x) 前先看符号。
\item \textbf{Tr 归一}：本文 $\Tr$ 取定义表示、长根平方 $=2$。若改成 $1/2$ 的迹，
(1.3) 的 $k/4\pi$ 要换成 $k/8\pi$（\src{354} 脚注 1 之后的说明）。不要在 (1.2) 上
补 $1/4$——原刊照印，无 $1/4$、无负号，见 \S\ref{app:errata}。
\item \textbf{页锚}：页边注“原文 p.$N$”指 CMP 印刷页；中译本 PDF 页 $=$ 印刷页 $+350$
只对\textbf{原文页图}成立，中译本本身是 40 页重排，不能拿译本 PDF 页套印刷页。
\end{enumerate}
